
\sourcenote{Genre: core rule (tier 1) --- an ordinary mathematical proof
at Putnam-exam-solution detail; the largest mechanized lemma set in this
book. The rule's mechanization is \emph{conditional}: the soundness
theorem is proved with a ledger of four named hypothesis families, of
which one is discharged outright below, two are reduced to
transit-certificate form, and one remains a hypothesis of the theorems
this section binds. The ledger remark near the end states each as
mathematics. Every hypothesis is used and remarked on; the formal-anchor
table maps each step to the Lean declaration that carries it.}

\textbf{Setting.}\enspace Fix a forward semiflow $\varphi$ on a metric
state space $X$ ($\Phi_0 = \mathrm{id}$,
$\Phi_{s+t} = \Phi_t \circ \Phi_s$), and a compact forward-invariant
arena $D \subseteq X$. Four standard notions, exactly as the mechanized
statements carry them:

\emph{(complete orbits)} $\gamma : \R \to X$ is a complete orbit when
$\gamma(s + t) = \Phi_t(\gamma(s))$ for all $s \in \R$, $t \geq 0$ ---
a two-sided trajectory compatible with the forward flow.

\emph{(maximal invariant part)} for $C \subseteq X$,
\[
\mathrm{Inv}(C) \;:=\; \{\, x \ :\ \exists\, \gamma \ \text{complete
orbit},\ \gamma(0) = x,\ \mathrm{range}\,\gamma \subseteq C \,\}.
\]
An \emph{invariant set} $S$ (the sense fit for a semiflow) is one that
is forward-invariant and satisfies $S \subseteq \mathrm{Inv}(S)$ ---
every point rides a complete orbit inside.

\emph{(limit sets)} $\omega(\gamma) := \bigcap_{T}
\mathrm{cl}\,\gamma([T, \infty))$ and $\alpha(\gamma) := \bigcap_{T}
\mathrm{cl}\,\gamma((-\infty, T])$.

\emph{(chains)} an $(\varepsilon, T)$-chain from $x$ to $y$ inside $D$
is a finite concatenation of orbit segments of duration $\geq T$ with
jumps of size $< \varepsilon$ at points of $D$ (generated by: a single
step $x, y \in D$, $t \geq T$,
$\mathrm{dist}(\Phi_t x, y) < \varepsilon$; and transitivity). A point
is \emph{chain-recurrent in $D$} when for every $\varepsilon > 0$ and
$T$ there is an $(\varepsilon, T)$-chain from $x$ to itself.

\textbf{The judgement.}\enspace
$\mathcal{M} \vDash_D \mathsf{Morse}(G, \rho)$ holds when some family
$\{M(v)\}_{v \in V(G)}$ --- indexed by the poset the graph $G$
presents by reachability --- satisfies:
each $M(v)$ compact and invariant; the family pairwise disjoint;
\emph{(location)} $M(v) \subseteq \rho(v)$;
\emph{(completeness of the cover)} every chain-recurrent point of
$\varphi|_D$ lies in some $M(v)$; and
\emph{(order)} every complete orbit in $D$ contained in no single
$M(v)$ has $\alpha(\gamma) \subseteq M(v)$ and
$\omega(\gamma) \subseteq M(w)$ for some $w < v$. The order runs
source-to-sink: $w < v$ means connecting orbits may run from $M(v)$
down to $M(w)$, so the $\alpha$-side vertex is the higher one. Empty
$M(v)$ are permitted --- the judgement is an ``at most this structure''
claim (recurrence confined to the $\rho(v)$, transitions confined to
the order), never an existence claim; the existence half is a
separately priced annotation tier (Conley indices, connection
matrices), deliberately absent here.

\textbf{The rule's data and the checker's recompute.}\enspace The
certificate supplies: the family object --- a finite lattice of
box-union blocks over one declared rational grid, each internal facet
carrying a strictly-inward sign row (the \rulename{QBlockInv} section's
row species; the same rows serve both rules); a cell map $h$ from the
lattice's cells to $V(G) \sqcup \{\mathsf{transit}\}$; per vertex cell
a containment row $T(b) \subseteq \rho(h(b))$; per transit cell a
polynomial $V$ with decrease rows $\mathcal{A}V + \varepsilon \leq 0$;
and well-formedness rows for $\rho$. The checker \emph{recomputes},
rather than trusts, the lattice readout: closure under union and
grid-level intersection, the join-irreducible members $b$ with unique
lower covers $b^-$, the cells $T(b) := \mathrm{cl}(b \smallsetminus
b^-)$, the induced cell order, convexity of $h$'s fibres and
monotonicity of $h$ on vertex cells, and acyclicity of $G$. Writing
$\mathrm{fib}(v) := \bigcup_{h(b) = v} T(b)$ for the fibre's
cell-union, the witness the rule assembles is
\[
M(v) \;:=\; \mathrm{Inv}\big(\mathrm{fib}(v)\big)
\]
--- the invariant part \emph{of the fibre's cell-union}, not the union
of per-cell invariant parts. The distinction is essential
(Remark~\ref{rem:qmd-mv}).

\begin{blemma}[$\mathrm{Inv}$ is located and invariant]
\label{lem:qmd-inv}
\emph{(i)} $\mathrm{Inv}(C) \subseteq C$.
\emph{(ii)} $\mathrm{Inv}(C)$ is an invariant set.
\end{blemma}
\begin{proof}
(i) The witnessing orbit passes through $x$ at time $0$ with range in
$C$. (ii) Two facts, both by time-shifting: shifting a complete orbit
by any $s$ ($u \mapsto \gamma(u + s)$) is again a complete orbit, and
every point $\gamma(s)$ of an orbit with range in $C$ therefore lies in
$\mathrm{Inv}(C)$ with the shifted orbit as witness. Forward
invariance: for $x \in \mathrm{Inv}(C)$ with orbit $\gamma$,
$\Phi_t x = \gamma(0 + t)$ is such a point. The complete-orbit-through
clause: the same statement at every shift shows
$\mathrm{range}\,\gamma \subseteq \mathrm{Inv}(C)$, so the witnessing
orbit itself runs inside $\mathrm{Inv}(C)$.
\end{proof}

\begin{blemma}[location]\label{lem:qmd-loc}
Under the containment rows, $\mathrm{fib}(v) \subseteq \rho(v)$, hence
$M(v) \subseteq \rho(v)$ by Lemma~\ref{lem:qmd-inv}(i).
\end{blemma}
\begin{proof}
A union of cells each inside $\rho(v)$ is inside $\rho(v)$.
\end{proof}

\begin{blemma}[forward confinement]\label{lem:qmd-conf}
If $x \in \mathrm{Inv}(C)$ then $\Phi_t x \in C$ for every
$t \geq 0$.
\end{blemma}
\begin{proof}
The witnessing orbit's forward half \emph{is} the forward orbit:
$\Phi_t x = \Phi_t(\gamma(0)) = \gamma(t) \in C$. This one-line lemma
is the section's quiet workhorse: a semiflow is deterministic
\emph{forward}, so membership in an $\mathrm{Inv}$ set constrains the
actual forward trajectory --- and nothing here assumes backward
uniqueness (a semiflow may branch backward; complete orbits through a
point need not be unique). The disjointness reduction below survives
precisely because it only ever uses forward halves.
\end{proof}

\begin{blemma}[transit certificates kill invariance]\label{lem:qmd-transit}
Call $C$ \emph{transit-certified} when there are $V : X \to \R$ and
$\varepsilon > 0$ with $V$ bounded below on $C$ and, along every
forward trajectory confined to $C$,
$V(\Phi_t x) \leq V(x) - \varepsilon\, t$. Then no point is
forward-confined to $C$; in particular
$\mathrm{Inv}(C) = \emptyset$.
\end{blemma}
\begin{proof}
Suppose $\Phi_t x \in C$ for all $t$, and let $m$ bound $V$ below on
$C$. By the Archimedean property choose $n$ with
$n\,\varepsilon > V(x) - m$; the decrease clause at horizon $n$ gives
$V(\Phi_n x) \leq V(x) - \varepsilon n < m$, contradicting
$\Phi_n x \in C$. The $\mathrm{Inv}$ statement follows by
Lemma~\ref{lem:qmd-conf}. (The certified rows are pointwise ---
$\mathcal{A}V + \varepsilon \leq 0$ on the cell, for the deterministic
generator $\mathcal{A} = f \cdot \nabla$ --- and reach the integrated
clause used here by integrating the derivative along the trajectory;
that bridge is ledger item G2 in Remark~\ref{rem:qmd-ledger}, stated
once and consumed by every certificate-shaped hypothesis in this
section.)
\end{proof}

\begin{blemma}[disjointness from overlap certificates]
\label{lem:qmd-disj}
If $C_1 \cap C_2$ is transit-certified, then
$\mathrm{Inv}(C_1)$ and $\mathrm{Inv}(C_2)$ are disjoint.
\end{blemma}
\begin{proof}
A common point is forward-confined to $C_1$ and to $C_2$
(Lemma~\ref{lem:qmd-conf} twice), hence to $C_1 \cap C_2$ --- which
Lemma~\ref{lem:qmd-transit} forbids. No backward uniqueness, no
injectivity of the flow maps.
\end{proof}

\begin{bremark}[why the certificate sits on the overlap ---
``strictness pays a second time'']\label{rem:qmd-overlap}
Distinct fibres' cell-unions intersect only in shared boundary faces
(cell interiors are disjoint by the readout's construction ---
generically facets, but possibly lower-dimensional: two cells may meet
in an edge or a corner alone), and on any certified facet incident to
the shared face the facet polynomial $\ell$ itself serves as the
overlap observable: the shared face lies in $\ell$'s zero set, the
certified nonzero normal component makes $\ell$ strictly monotone
along the flow near the facet, while confinement to the face would
force $\ell \equiv 0$.
The strict margin of the \rulename{QBlockInv} rows is thus consumed a
second time, as a transit certificate for facet regions. One might
hope to weaken the hypothesis to per-piece certificates --- each
$C_i$ separately transit-certified --- and conclude something about
the union; that is false. On the circle rotation, two overlapping
closed arcs each carry a monotone angular observable (a transit
certificate for the arc), yet the union is the whole circle, every
point recurrent: certificates on pieces do not compose across an
overlapping union, because no global observable exists. The overlap
form is the correct, and certifiable, hypothesis.
\end{bremark}

\begin{blemma}[compactness of the invariant parts]\label{lem:qmd-cpt}
Suppose $X$ is Hausdorff, each flow map $\Phi_t : X \to X$ is
continuous, and $C$ is compact. Then $\mathrm{Inv}(C)$ is compact.
Consequently each $\mathrm{fib}(v)$, a finite union of compact cells,
has compact invariant part.
\end{blemma}
\begin{proof}
Let $K \subseteq X^{\R}$ (product topology) be the set of complete
orbits with range in $C$:
\[
K \;=\; \Big(\bigcap_{s \in \R,\ t \geq 0}
\{\gamma : \gamma(s+t) = \Phi_t(\gamma(s))\}\Big)
\;\cap\; \prod_{r \in \R} C .
\]
Each orbit constraint is the equalizer of two continuous maps ---
evaluation at $s + t$, and $\Phi_t \circ (\text{evaluation at } s)$
--- hence closed, $X$ being Hausdorff; the product
$\prod_{\R} C$ is compact by Tychonoff (the one genuinely infinitary
step of this section); a closed subset of a compact set is compact. And
$\mathrm{Inv}(C)$ is exactly the image of $K$ under evaluation at $0$,
a continuous map, hence compact. The fibre statement: a finite union
of compact cells is compact, and the above applies to it.
\end{proof}

\begin{blemma}[the cover assembles from the cell level]
\label{lem:qmd-cover}
Suppose \emph{(cell-level cover)} every chain-recurrent point of
$\varphi|_D$ lies in $\mathrm{Inv}(T(b))$ for some cell $b$, and
\emph{(transit rows)} every transit cell is transit-certified. Then
every chain-recurrent point of $\varphi|_D$ lies in some
$M(v) = \mathrm{Inv}(\mathrm{fib}(v))$.
\end{blemma}
\begin{proof}
Let $x$ be chain-recurrent, $x \in \mathrm{Inv}(T(b))$. If $h(b) =
\mathsf{transit}$, Lemma~\ref{lem:qmd-transit} makes
$\mathrm{Inv}(T(b))$ empty --- contradiction; so $h(b)$ is a vertex
$v$. Then $T(b) \subseteq \mathrm{fib}(v)$, and $\mathrm{Inv}$ is
monotone in its region (a complete orbit inside the cell is inside the
fibre's union), so $x \in \mathrm{Inv}(\mathrm{fib}(v))$.
\end{proof}

\begin{bthm}[soundness of \rulename{QMorseDischarge}, at the
discharged floor]\label{thm:qmd-sound}
Suppose:
\emph{(model class)} $X$ is a Hausdorff metric space, each flow map
$\Phi_t$ is continuous, and every cell is compact;
\emph{(containment rows)} $T(b) \subseteq \rho(h(b))$ for every vertex
cell;
\emph{(overlap certificates)} for $v \neq w$, the overlap
$\mathrm{fib}(v) \cap \mathrm{fib}(w)$ is transit-certified;
\emph{(cell-level cover)} every chain-recurrent point of $\varphi|_D$
lies in some cell's invariant part;
\emph{(transit rows)} every transit cell is transit-certified; and
\emph{(order clause)} every complete orbit in $D$ contained in no
single $M(v)$ connects a strictly higher fibre's invariant part to a
strictly lower one, in the sense of the judgement's order clause with
$M(v) = \mathrm{Inv}(\mathrm{fib}(v))$.
Then $\mathcal{M} \vDash_D \mathsf{Morse}(G, \rho)$.
\end{bthm}
\begin{proof}
Take the family $M(v) := \mathrm{Inv}(\mathrm{fib}(v))$ and check the
witness clauses. \emph{Location}: Lemma~\ref{lem:qmd-loc}.
\emph{Invariance}: Lemma~\ref{lem:qmd-inv}(ii). \emph{Compactness}:
Lemma~\ref{lem:qmd-cpt} under the model-class facts.
\emph{Disjointness}: Lemma~\ref{lem:qmd-disj} applied to the overlap
certificates. \emph{Cover}: Lemma~\ref{lem:qmd-cover} from the
cell-level cover and the transit rows. \emph{Order}: the hypothesis,
verbatim. All six clauses hold, so the family witnesses the judgement.
\end{proof}

\begin{bremark}[the gap ledger, as mathematics]\label{rem:qmd-ledger}
The rule was first mechanized as the bare witness assembly: location
and invariance proved outright, and four clauses --- labelled G3--G6 in
the development, the labels this book adopts as the rule's hypothesis
names --- taken as opaque interface hypotheses. The discharge stratum
bound by this section moves the ledger to the state
Theorem~\ref{thm:qmd-sound} exhibits:

\emph{G3 (compactness of the invariant parts): proven.}
Lemma~\ref{lem:qmd-cpt} derives it from model-class facts alone ---
Hausdorff state space, spacewise-continuous flow, compact cells. No
compactness hypothesis survives at the fibre level.

\emph{G4 (disjointness of distinct fibres' invariant parts): reduced.}
Lemma~\ref{lem:qmd-disj} trades it for transit certificates on
pairwise overlaps; the residue is the overlap observable of
Remark~\ref{rem:qmd-overlap} --- shared-facet functionals whose
certified strict signs are the raw material --- together with the G2
bridge below.

\emph{G5 (completeness of the chain-recurrence cover): reduced.}
Lemma~\ref{lem:qmd-cover} trades the fibre-level cover for the
cell-level one. The residue is the margin argument: a chain-recurrent
point's chains stay within its own cell, because the strict facet
margin is a uniform bound on admissible chain jumps near facets --- an
$\varepsilon$-chain with $\varepsilon$ below the margin can no more
exit a block than an orbit can. The classical home of this
containment is Conley's theory \cite{conley1978}; its algorithmic
form is the chain-recurrence algorithm of
Kalies--Mischaikow--VanderVorst \cite{kmv2005}.

\emph{G6 (the order clause): remains a hypothesis} of the theorems
bound in this section. Its mathematical content is the orbit
combinatorics of the block lattice: the \emph{eventual member} of a
forward orbit --- the minimal lattice member containing a forward
tail, well-defined because the tail-containing family is finite,
nonempty, and closed under the lattice meet --- descends the cell
order and stabilizes, giving every complete orbit in $D$ extreme
cells $b_\omega \leq b_\alpha$ that bracket every visited cell; the
fibre-convexity case analysis then converts bracketing into the
witness's order clause (the extremes are not transit, by
Lemma~\ref{lem:qmd-transit} against their nonempty invariant limit
sets; a same-fibre pair traps the whole orbit in its fibre by
convexity, contradicting the ``no single $M(v)$'' hypothesis; distinct
fibres yield the strict order by monotonicity of $h$). Successor
modules of the accompanying development mechanize exactly this
assembly --- the cell order entering as theorem-interface data, the
eventual-member construction proved against combinatorial row
hypotheses --- and further reduce the G4/G5 residues to per-facet
crossing and mean-value certificates; those modules are beyond what
this section binds, and the rule's conditional badge is judged against
the bound statements.

\emph{G2 (the pointwise-to-integrated bridge), named for
completeness:} every ``transit-certified'' hypothesis above is the
integrated semantic shadow of pointwise rows
$\mathcal{A}V + \varepsilon \leq 0$; the bridge is one-variable
calculus along trajectories, and its surviving analytic content is
the identification of the certified polynomial rate with the
trajectory-wise derivative on the region --- the same model-class
profile warrant as ledger item (a) of the \rulename{QBlockInv}
section, whose G1 (block invariance) completes the shared ledger.
\end{bremark}

\begin{bremark}[the witness is $\mathrm{Inv}$ of the union, not the
union of $\mathrm{Inv}$s]\label{rem:qmd-mv}
Define instead $M'(v) := \bigcup_{h(b) = v} \mathrm{Inv}(T(b))$ and
the order clause fails: a complete orbit crossing between two cells of
the \emph{same} fibre lies in no per-cell invariant part (its range is
in neither cell alone), so it belongs to no $M'(v)$ while also
connecting nothing strictly ordered --- the ``orbit contained in no
Morse set'' case triggers with no higher/lower pair to name. The
correct definition absorbs such orbits into their fibre's $M(v)$,
which is exactly the role of fibre convexity: every cell between the
orbit's extremes maps to the same vertex, so the whole orbit lies in
the fibre's union and hence in its invariant part. Without convexity
the judgement could assert a strict order the certificate does not
enforce --- convexity of $h$'s fibres is well-formedness, not a
nicety, and the checker recomputes it.
\end{bremark}

\begin{bremark}[relation to the literature]
The judgement is Conley's Morse decomposition \cite{conley1978} in
certified form: a decomposition indexed by a two-element order is an
attractor--repeller pair, and the general form is the lattice-refined
iteration of that idea; the completeness clause is the containment of
chain recurrence in the Morse sets, and the certified staircase of
lattice members is a finite shard of a complete Lyapunov function
(Conley's fundamental theorem). The block-lattice route --- attracting
blocks closed under union and intersection, join-irreducibles, the
tessellation into cells $T(b)$, and Birkhoff duality
\cite{birkhoff1937} between the attractor lattice and the Morse-graph
poset --- is the program of Kalies--Mischaikow--VanderVorst's lattice
series \cite{kmv2016,kmv2022}; the chain-recurrence cover's
algorithmic form is their earlier chain-recurrence algorithm
\cite{kmv2005}; and the computational tradition of certifying global
dynamics by index machinery is surveyed in Mischaikow--Mrozek
\cite{mischaikowmrozek2002}. What the rule adds to that tradition is
the artefact discipline: those pipelines end in trusted computation,
while here every analytic claim rides certified sign rows --- base
kernel sequents in exact $\Q$ --- plus combinatorics the checker
recomputes, with the semantic warrant being the lemma chain above.
\end{bremark}

\begin{anchors}
witness assembly, location + invariance proved & \texttt{qMorseDischarge\_sound} (QtslSoundness.lean); \texttt{maxInv\_subset}, \texttt{maxInv\_isInvariant}\\
Lemma \ref{lem:qmd-inv} & \texttt{maxInv\_subset}, \texttt{maxInv\_invariant}, \texttt{maxInv\_isInvariant}\\
Lemma \ref{lem:qmd-loc} & \texttt{fiberUnion\_subset}; rows \texttt{hContain}\\
Lemma \ref{lem:qmd-conf} & \texttt{maxInv\_forward\_confined} (Discharge.lean)\\
Lemma \ref{lem:qmd-transit} & \texttt{HasTransitCert}, \texttt{HasTransitCert.not\_forward\_confined}, \texttt{HasTransitCert.maxInv\_empty}; the bare form \texttt{transit\_maxInv\_empty} (QtslSoundness.lean)\\
Lemma \ref{lem:qmd-disj} & \texttt{HasTransitCert.disjoint\_maxInv}; hypothesis \texttt{hOverlap}\\
Lemma \ref{lem:qmd-cpt} & \texttt{maxInv\_isCompact} (Tychonoff route), \texttt{fiberUnion\_isCompact}\\
Lemma \ref{lem:qmd-cover} & \texttt{crCover\_assemble}; \texttt{maxInv\_mono}, \texttt{subset\_fiberUnion}; hypotheses \texttt{hCell}, \texttt{hTransitCert}\\
Theorem \ref{thm:qmd-sound} & \texttt{qMorseDischarge\_sound\_modelClass}; order hypothesis \texttt{hOrder} (ledger G6)\\
judgement semantics & \texttt{MorseWitness}, \texttt{MorseSat}, \texttt{EpsChainIn}, \texttt{ChainRecurrentIn}, \texttt{omegaSet}, \texttt{alphaSet}\\
G6's successor mechanization (not bound here) & \texttt{orderClause\_assemble}, \texttt{eventualMember\_interval}, \texttt{qMorseDischarge\_sound\_orderAssembled}, \texttt{\_eventualMember} (development)\\
\end{anchors}
